\addtocontents{toc}{\protect\newpage}
\section{\MakeUppercase{Computational Methods}} 
\label{sec:Deepdiffusion}

\counterwithin{figure}{section}
\counterwithin{table}{section}


A few computational techniques exist for transferring quantities from between cells.
These methods -- Finite Difference, Finite Volume, and Smooth Particle Hydrodynamics -- excel at capturing different aspects of fluid flow.

A supernova expanding into a dense medium will generate an extremely strong shock.
The Finite Volume method is particularly well-suited for this situation because it is designed to conserve the flux of quantities between cells via the Riemann problem.
Since the grid is discretized, neighboring cells contain different quantities. The difference in a quantity like density may be relatively small, but the interface between these cells produces discontinuity.
The way this system evolves in time and quantities are exchanged between the cells can be solved using the Riemann problem in a conserved manner.
Thus, the Finite Volume method excels at handling shock waves formed in supersonic flows while retaining numerical stability since the problem at hand is built into the method.